\subsection{The One-Dimensional Pairing Experiment}
Take $H\sim\operatorname{Unif}[1,4]$ and
$R\sim\operatorname{Unif}[.01,1]$, with equal versus opposite quantile
pairing and independent fair noise signs. Both full separate distributions
and \eqref{eq:matched-kernel} agree. Here noise magnitude is a fixed function
of curvature, so the joint distribution lies on curves rather than having
a smooth density in both variables. Implementation uses
uniform 52-bit midpoint integers and the same rounded affine map at $U$
and $1-U$, making implemented Hessian marginals exactly equal.

Use $\alpha=.1$, $\beta=.2$, $\epsilon=.03$, $\lambda=.1$ and three
fixed learning rates. Each arm has 256 paired runs, starts $0,\pm8\sqrt\eta$
and three 17-sample windows at rescaled times $16{:}32$, $48{:}64$,
$112{:}128$. An independent calculation made before these simulations gives
a joint 99\% interval $[.472216,.633164]$ for the limiting
opposite-rank/equal-rank AdamW error ratio. Table~\ref{tab:scalar-errors} uses the zero start
and final observation window:
\begin{table}[htbp]
\caption{Supporting One-Dimensional Pairing Comparison}\label{tab:scalar-errors}
\centering

\begin{tabular}{ccc}
\toprule
$\eta$ & Opposite/matching error & Approx. pointwise 95\% interval\\
\midrule
$5\cdot10^{-4}$ & $.550718$ & $[.540987,.560802]$\\
$2\cdot10^{-4}$ & $.543005$ & $[.533611,.552703]$\\
$8\cdot10^{-5}$ & $.544562$ & $[.534866,.554215]$\\
\bottomrule
\end{tabular}
\end{table}
Figure~\ref{fig:scalar-couplings} also shows the absolute normalized errors
and the independent prediction intervals. All 27 paired AdamW differences
have negative upper endpoints in the planned approximate simultaneous band.
All 54 comparisons with exact SGD errors have intervals containing zero.
Both AdamW arms have higher absolute error than SGD at these learning rates.

The independent calculation gives point values $k_{\rm equal}=1.965666$,
$k_{\rm opposite}=3.583588$ and $D=.781407$. They predict common-step
risk ratio $.548519$, error ratio $.300873$ at matched local return speed,
and local return-speed ratio $3.323661$ at equal leading error.
The respective efficiency quantities
$D/k^2$ are $.202236$, $.0608473$, and $.0498077$ for equal AdamW,
opposite AdamW, and SGD. Lower values mean less error at the same local
return speed. The pairing effect remains, and SGD still has the lowest
value in this example.

\subsection{An Exploratory Check of Correlations Over Time}
The limiting stationary process and fourth-moment bounds imply
\[
 \lim_{\eta\downarrow0}
 \frac{\E[u_0u_{\lfloor t/\eta\rfloor}]}\eta
       =\frac D{2k}e^{-k|t|}.
\]
A check chosen after examining the data uses all saved adjacent-time pairs
and every setting, start and time window, with the existing 10,000 recorded
resamples of whole runs. It draws no new values and fits no drift. The
uncentered mean product at the two times is divided by their average squared parameter value;
the independent prediction uses exact integer time differences.
All 108 approximate simultaneous bands intersect the independently predicted
limiting intervals. At all 27 AdamW learning-rate/start/window combinations,
the equal-rank arm's simultaneous lower endpoint exceeds the opposite-rank
arm's upper endpoint. These quantities are correlated. The check is
exploratory and gives no prospective test of time to stationarity.

\subsection{Earlier Disagreements and a More Direct Response Test}
A separate earlier rare-pulse study let the denominator offset tend to zero.
The main theorem here instead keeps it fixed and positive. The earlier grid
missed the leading predictions: only 15 of 108 planned simultaneous moment
bands intersected the independent prediction bounds. After examining those
data, we derived a first correction and fixed a new test at smaller learning
rates before generating its trajectories. Its 36 correlated point estimates
lie inside the corrected prediction intervals, but all 36 simultaneous bands
intersect both the leading and corrected intervals. Unknown finite-step
errors and different observation clocks prevent a decisive distinction.
The companion archive keeps those results and proofs; none is needed for
Theorem~\ref{thm:main}.

A new experiment compares two weight decays while holding the gradient
sampler and observation clock fixed. Subtracting the two results cancels the
shared leading prediction. All 18 planned paired differences have negative
upper endpoints in the approximate simultaneous band, as predicted by the
independently calculated first correction. Their normalized values are
compatible with its prediction intervals. Appendix~\ref{app:empirical}
gives both rates and the limits of this comparison. This stronger response
evidence leaves the earlier disagreements and the unknown finite-step
error bounds unresolved.

